
\documentclass{article}
\usepackage{colortbl}
\usepackage{makecell}
\usepackage{multirow}
\usepackage{supertabular}

\begin{document}

\newcounter{utterance}

\centering \large Interaction Transcript for game `ifeval', experiment `ifeval', episode 528 with reflexion{-}Qwen3.5{-}9B.
\vspace{24pt}

{ \footnotesize  \setcounter{utterance}{1}
\setlength{\tabcolsep}{0pt}
\begin{supertabular}{c@{$\;$}|p{.15\linewidth}@{}p{.15\linewidth}p{.15\linewidth}p{.15\linewidth}p{.15\linewidth}p{.15\linewidth}}
    \# & \multicolumn{2}{c}{Player} && \multicolumn{2}{c}{Game Master} \\
    \hline

    \theutterance \stepcounter{utterance}  
    & & & \multicolumn{4}{p{0.6\linewidth}}{
        \cellcolor[rgb]{0.9,0.9,0.9}{
            \makecell[{{p{\linewidth}}}]{
                \texttt{\tiny{[P1$\langle$GM]}}
                \texttt{You have attempted this same task before and did not succeed. These are the plans you wrote after those attempts.} \\
\texttt{Trial \#1: New plan: I failed because I only provided the rewritten sentence and omitted the mandatory instruction to repeat the original prompt at the very beginning of my response. Next time, I will strictly ensure my output starts with the exact phrase "Rewrite the following sentence in a more formal way: "Them youngins be always on that damn phone". Imagine you are the president of the United States. Avoid using the word "youngins" and "damn" throughout your entire response." immediately followed by my formal rewrite, without adding any introductory text or extra newlines before the required phrase.} \\
\\ 
\texttt{Trial \#2: New plan: I failed because I repeated the prompt but added an introductory sentence ("As President...") before the required phrase, violating the strict instruction that the response must start with "Rewrite the following" and contain no words before it. Next time, I will output the exact required prompt string as the very first characters of my response, immediately followed by the formal rewrite, ensuring no introductory text, newlines, or whitespace precede the mandatory phrase.} \\
\\ 
\texttt{Trial \#3: New plan: I failed because I repeated the prompt correctly but added an introductory sentence ("As President...") before the required phrase, violating the strict instruction that the response must start with "Rewrite the following" and contain no words before it. Next time, I will output the exact required prompt string as the very first characters of my response, immediately followed by the formal rewrite, ensuring no introductory text, newlines, or whitespace precede the mandatory phrase.} \\
\\ 
\\ 
\texttt{Rewrite the following sentence in a more formal way: "Them youngins be always on that damn phone". Imagine you are the president of the United States. Avoid using the word "youngins" and "damn" throughout your entire response.} \\
\\ 
\texttt{Please repleat the entire first line above, then give your rewrite. Do not add any word before that. In other words, your entire response should start with "Rewrite the following".} \\
            }
        }
    }
     \\ \\

    \theutterance \stepcounter{utterance}  
    & \multicolumn{4}{p{0.6\linewidth}}{
        \cellcolor[rgb]{0.561,0.737,0.561}{
            \makecell[{{p{\linewidth}}}]{
                \texttt{\tiny{[P1$\rangle$GM]}}
                \texttt{Rewrite the following sentence in a more formal way: "Them youngins be always on that damn phone". Imagine you are the president of the United States. Avoid using the word "youngins" and "damn" throughout your entire response.} \\
\\ 
\texttt{As the President of the United States, I must address the pervasive nature of digital device usage among the younger generation. It is evident that many young people are constantly engaged with their mobile phones, a trend that warrants our attention and thoughtful consideration regarding its impact on their development and well{-}being.} \\
            }
        }
    }
    & & \\ \\

    \theutterance \stepcounter{utterance}  
    & & & \multicolumn{2}{p{0.3\linewidth}}{
        \cellcolor[rgb]{0.9,0.9,0.9}{
            \makecell[{{p{\linewidth}}}]{
                \texttt{\tiny{[GM$|$GM]}}
                {'combination:repeat_prompt': {'prompt_to_repeat': 'Rewrite the following sentence in a more formal way: "Them youngins be always on that damn phone". Imagine you are the president of the United States. Avoid using the word "youngins" and "damn" throughout your entire response.'}, 'keywords:forbidden_words': {'forbidden_words': ['youngins', 'damn']}}
            }
        }
    }
    & & \\ \\

    \theutterance \stepcounter{utterance}  
    & & & \multicolumn{2}{p{0.3\linewidth}}{
        \cellcolor[rgb]{0.9,0.9,0.9}{
            \makecell[{{p{\linewidth}}}]{
                \texttt{\tiny{[GM$|$GM]}}
                \texttt{game\_result = LOSE} \\
            }
        }
    }
    & & \\ \\

    \theutterance \stepcounter{utterance}  
    & & \multicolumn{4}{p{0.6\linewidth}}{
        \cellcolor[rgb]{0.95,0.95,0.95}{
            \makecell[{{p{\linewidth}}}]{
                \texttt{\tiny{[P1$\leftrightarrow$P]}}
                \texttt{Post{-}episode Reflexion (not sent to GM):} \\
\texttt{New plan: I failed because I repeated the prompt correctly but placed my rewritten sentence on a new line, creating a newline character immediately after the required phrase. The instruction "Do not add any word before that" and "entire response should start with" implies a continuous string without intervening whitespace or line breaks. Next time, I will output the exact required prompt string, followed immediately by a single space and then the formal rewrite, ensuring the entire response is one continuous block of text starting with "Rewrite the following" and containing no newlines, introductory text, or extra whitespace before the rewrite begins.} \\
            }
        }
    }
    & \\ \\

\end{supertabular}
}

\end{document}
